\documentclass[10pt,a4paper]{amsart}
\usepackage[margin=19mm]{geometry}
\usepackage{amsmath,amssymb,mathtools}
\usepackage[scr=rsfs,cal=euler]{mathalfa}
\usepackage{xfrac}
\usepackage[unicode,hidelinks,bookmarks=false]{hyperref}
\newcommand{\ie}{i.e.\ }
\newcommand{\cf}{cf.\ }
\newcommand{\resp}{respectively\ }
\newcommand{\Subsection}{\subsection}
\providecommand{\nobreakdash}{\nobreak\hbox{-}}
\setlength{\emergencystretch}{2em}
\multlinegap0pt
\usepackage{bm}
\usepackage{bbm}
\usepackage{dsfont}
\usepackage{xspace}
\usepackage{cancel}
\usepackage{mathtools}
\usepackage[shortlabels]{enumitem}
\usepackage{amsthm}
\makeatletter
\@ifundefined{thm}{\newtheorem{thm}{Thm}}{}
\@ifundefined{lem}{\newtheorem{lem}[thm]{Lemma}}{}
\@ifundefined{notation}{\newtheorem{notation}[thm]{Notation}}{}
\makeatother
\newcommand{\id}{\operatorname{id}}
\newcommand{\ootimes}{\, \overline{\otimes} \,}
\title[Exotic full factors via weakly coarse bimodules]{Exotic full factors via weakly coarse bimodules}
\author{David Gao, David Jekel, Srivatsav Kunnawalkam Elayavalli, Gregory Patchell}
\date{Source: arXiv:2602.00930v1 (CC BY 4.0)}
\begin{document}
\maketitle

\noindent\emph{Source and licence.} Exotic full factors via weakly coarse bimodules. Version arXiv:2602.00930v1 (CC BY 4.0), distributed under the Creative Commons Attribution 4.0 International licence, as recorded in the arXiv licence metadata for this exact version and shown on the abstract page https://arxiv.org/abs/2602.00930v1. Full rights notice: licenses/arxiv-2602.00930-CC-BY-4.0.txt.\par\smallskip

\noindent\textit{Editorial scope.} Counted unit: the Lem whose source label is \texttt{lem: word decomposition 3} in the article (source0.tex lines 502--507, source label \texttt{lem: word decomposition 3}), delivered together with the article's own complete original proof of it (source0.tex lines 509--553, one \texttt{proof} environment of 45 lines). No mathematics is written, repaired or re-derived: statement and proof are the article's own text line for line. Required context retained uncounted: not: free product setup (lines 368--375); lem: word decomposition (lines 389--397); lem: word decomposition 2 (lines 450--462). Every definition, display and label of those lines is retained unchanged. Omitted from this package and not redistributed: 1--367, 376--388, 398--449, 463--501, 508--508, 554--650. Nothing inside a retained excerpt is omitted or altered: every retained body is delivered exactly as the article's lines are written, so no omission receipt of any kind accompanies this package. Citations: the retained excerpts carry no citation command at all, so no bibliography entry of the article is transcribed and no reference is redistributed. Reference adaptation: none was needed and none was made; every reference inside the delivered unit resolves inside it, so no unresolved reference or citation is produced. Printed slips kept literal: no printed word, symbol, hypothesis or number of the counted statement or of its proof was altered. Layout and class: the article's own class is \texttt{amsart} with packages including amsmath, amsthm, amssymb, eucal, cite, upgreek, hyperref, enumerate, mathrsfs, blindtext, scrextend, bm, color, easylist; the wrapper is the lane's pinned \texttt{amsart} editorial wrapper at 10pt on A4, which restates the theorem-like environments the retained text needs and adds the article's own macro definitions that the retained text needs (\texttt{\textbackslash id}, \texttt{\textbackslash ootimes}), with extra wrapper packages (bm, bbm, dsfont, xspace, cancel, mathtools, enumitem). The delivered numbering is the unit's own. Assets: no retained span contains a figure environment, a graphics-inclusion command, a PDF-page inclusion, a table or a TikZ picture; no image file is copied into this package, the package declares no asset dependency and no image file is redistributed with it. Licence: version arXiv:2602.00930v1 of this article is distributed under the Creative Commons Attribution 4.0 International licence, as recorded in the arXiv licence metadata for this exact version and shown on the abstract page https://arxiv.org/abs/2602.00930v1. A full rights notice is delivered at licenses/arxiv-2602.00930-CC-BY-4.0.txt. Characters: every retained excerpt is pure ASCII, so the delivered engine, XeLaTeX with its default Latin Modern text font, reports no missing character.
\begin{notation} \label{not: free product setup}
	Let $n \geq 1$.  For $j=1, \dots, n$, let $(A,\varphi)$ be a von Neumann algebra with faithful normal state and $B_j \subset B$ a von Neumann subalgebra with state-preserving conditional expectation $E_{B_j}$; let $\varphi_j = \varphi|_{B_j}$.  Let $(C_j,\psi_j)$ be another von Neumann algebra with faithful normal state.  Define von Neumann algebras $M_0 \subset M_1 \subset \dots \subset M_n$ with states $\omega_j$ inductively by
	\begin{align*}
		(M_0,\omega_0) &= (A,\varphi) \\
		(M_{j+1},\omega_{j+1}) &= (M_j,\omega_j) *_{(B_j,\varphi_j)} [(B_j,\varphi_j) \ootimes (C_j,\psi_j)],
	\end{align*}
	where $(B_j,\varphi_j) \subset (A,\varphi) = (M_0,\omega_0)$ is regarded as a subalgebra of $(M_j,\omega_j)$ in the natural way, and we also use the natural inclusion of $(B_j,\varphi_j)$ into $(B_j,\varphi_j) \ootimes (C_j,\psi_j)$ with the conditional expectation given by $\id \otimes \psi_j$.  We also regard $(A,\varphi)$, $(B_j,\varphi_j)$, and $(C_j,\psi_j)$ as subalgebras of $(M_n,\omega_n)$ in the natural way.
\end{notation}

\begin{lem} \label{lem: word decomposition}
	Consider the setup of Notation \ref{not: free product setup}.  For each alternating word $w = j_1 \dots j_\ell$ on the alphabet $\{0,\dots,n\}$, let $K_w$ be the span of products of the form $x_1 \dots x_\ell$, where
	\begin{enumerate}[(a)]
		\item if $j_i > 0$, then $x_i \in C_{j_i}$ with $\psi_{j_i}(x_i) = 0$;
		\item if $j_i = 0$, then $x_i \in A$;
		\item if $j_i = 0$ and $i < \ell$, then $E_{B_{j_{i+1}}}[x_i] = 0$.
	\end{enumerate}
	Then the $K_w$'s span a weakly dense subset of $M_n$.
\end{lem}

\begin{lem} \label{lem: word decomposition 2}
	Consider the setup of Notation \ref{not: free product setup}, and let $K_w$ be as in Lemma \ref{lem: word decomposition}.  Suppose $w = j_1 \dots j_\ell$ is an alternating word on $\{0,\dots,n\}$.  Consider a product $x_1 \dots x_\ell$ satisfying
	\begin{itemize}
		\item[\emph{(a)}] if $j_i > 0$, then $x_i \in C_{j_i}$ with $\psi_{j_i}(x_i) = 0$;
		\item[\emph{(b)}] if $j_i = 0$, then $x_i \in A$;
		\item[\emph{(c')}] if $1 < i < \ell$ with $j_i = 0$ and $j_{i+1} = j_{i-1}$, then $E_{B_{j_{i-1}}}[x_i] = 0$.
	\end{itemize}
	Then
	\begin{enumerate}[(1)]
		\item If the word $w$ contains the letter $n$ at least once, then $E_{M_{n-1}}[x_1 \dots x_\ell] = 0$.
		\item Moreover, if the word $w$ contains any nonzero letter, then $E_A[x_1 \dots x_\ell] = 0$.
	\end{enumerate}
\end{lem}

\begin{lem} \label{lem: word decomposition 3}
Regard $C$ as a von Neumann subalgebra of $(M_n,\omega_n)$.  Let $H_0$ be the span of $K_w$ for alternating words $w$ on $\{0,\dots,n\}$ which start with the letter $0$.  Then for $\xi$, $\eta \in H_0$ and $x, y \in C$, we have
\begin{equation} \label{eq: inner product equality}
\omega_n((x\xi)^*(y\eta)) = \psi(x^*y) \omega_n(\xi^* \eta).
\end{equation}
\end{lem}

\begin{proof}
By linearity, it suffices to consider the following case:  Let $w_0 = j_1 \dots j_\ell$ and $w_0' = j_1' \dots j_{\ell'}'$ be words in $\{0,\dots,n\}$ that start with zero, and suppose that $\xi = \xi_1 \dots \xi_\ell$ and $\eta = \eta_1 \dots \eta_{\ell'}$ are products satisfying the conditions (a)-(c) of Lemma \ref{lem: word decomposition} with respect to $w_0$ and $w_0'$, respectively.  Let $w_1 = k_1 \dots k_m$ and $w_1' = k_1' \dots k_{m'}'$ be alternating words on $\{1,\dots,n\}$ and suppose that $x = x_1 \dots x_m$ and $y = y_1 \dots y_{m'}$ are products of elements in $\ker(\psi)$ with $x_i \in C_{k_i}$ and $y_i \in C_{k_i'}$.

Note that the products $x_1 \dots x_m$ and $y_1 \dots y_{m'}$ also satisfy conditions (a)-(c) of Lemma \ref{lem: word decomposition}.  In fact, the products $x_1 \dots x_m \xi_1 \dots \xi_\ell$ and $y_1 \dots y_{m'} \eta_1 \dots \eta_{\ell'}$ satisfy (a)-(c) as well.  Indeed, since $j_1 = j_1' = 0$, the words are alternating.  Moreover, since all the occurrences of $0$ are in the right part of the word, we see that conditions (b) and (c) are satisfied.

The proof of the claim proceeds by induction on $m$ and $m'$.

\textbf{Case 1:} Suppose that $m$ and $m'$ are both zero.  Then $x$ and $y$ are empty products, which by convention are $1$.  Hence, \eqref{eq: inner product equality} is trivial. 

\textbf{Case 2:} Suppose that $m = 0$ and $m' > 0$.  Since $w_0$ and $w_0'$ start with $0$ and $w_1'$ does not contain $0$, the word $w_0^* w_1' w_0' = j_{\ell} \dots j_1 k_1' \dots k_{m'}' j_1' \dots j_{\ell'}'$ is alternating.  Moreover, the product $\xi_\ell^* \dots \xi_1^* y_1 \dots y_{m'} \eta_1 \dots \eta_{\ell'}$ satisfies (a), (b), (c') of Lemma \ref{lem: word decomposition 2} for this word, using the fact that $\xi_1 \dots \xi_\ell$ and $y_1 \dots y_{m'} \eta_1 \dots \eta_{\ell'}$ satisfy (a)-(c).  Hence, by Lemma \ref{lem: word decomposition 2}, we have
\[
\omega_n(\xi^* y \eta) = \omega_n(\xi_\ell^* \dots \xi_1^* y_1 \dots y_{m'} \eta_1 \dots \eta_{\ell'}) = 0.
\]
On the other hand, by free independence of the $C_j$'s, $\psi(y) = 0$.  Thus, both sides of \eqref{eq: inner product equality} are zero.

\textbf{Case 3:} Suppose that $m > 0$ and $m' = 0$.  The argument is symmetrical to Case 2.

\textbf{Case 4:} Suppose that $m > 0$ and $m' > 0$ and $k_1 \neq k_1'$.  Then the word $w_0^* w_1^* w_1' w_0'$ is alternating and the product $\xi_\ell^* \dots \xi_1 x_m^* \dots x_1^* y_1 \dots y_{m'} \eta_1 \dots \eta_{\ell'}$ satisfies conditions (a), (b), (c') of Lemma \ref{lem: word decomposition 2}. Therefore, $\omega_n(\xi^*x^*y\eta) = 0$.  On the other hand, since $w_1^*w_1'$ is an alternating word on $\{1,\dots,n\}$, we have $\psi(x^*y) = 0$ using free independence of $C_1$, \dots, $C_n$.  Hence, both sides of \eqref{eq: inner product equality} are zero.

\textbf{Case 5:} Suppose that $m > 0$ and $m' > 0$ and $k_1 = k_1'$.  Then the word
\[
j_\ell \dots j_1 k_m \dots k_1 k_2' \dots k_{m'}' j_1' \dots j_{\ell'}'
\]
is alternating and the product
\[
\xi_\ell^* \dots \xi_1^* x_m^* \dots x_2^* (x_1^*y_1 - \psi(x_1^*y_1)) y_2 \dots y_{m'} \eta_1 \dots \eta_{\ell'}
\]
satisfies conditions (a), (b), (c') of Lemma \ref{lem: word decomposition 2}.  Therefore,
\[
\omega_n(\xi^* x_m^* \dots x_2^* (x_1^*y_1 - \psi(x_1^*y_1)) y_1 \dots y_{m'} \eta) = 0,
\]
or
\[
\omega_n(\xi^* x^* y \eta) = \psi(x_1^*y_1) \omega_n(\xi^* x_m^* \dots x_2^* y_2 \dots y_{m'} \eta).
\]
Using the induction hypothesis,
\[
\omega_n(\xi^* x^* y \eta) = \psi(x_1^*y_1) \psi(x_m^* \dots x_2^* y_2 \dots y_{m'}) \omega_n(\xi^*\eta).
\]
By applying the same equality in the case that $\xi = \eta = 1$ (or using free independence),
\[
\psi(x^*y) = \psi(x_1^*y_1) \psi(x_m^* \dots x_2^* y_2 \dots y_{m'}),
\]
and thus, we obtain \eqref{eq: inner product equality} as desired.
\end{proof}

\end{document}
